\section*{الفصل \ref{ch:g11:trigo} --- حساب المثلثات: دائرة الوحدة}

\begin{solution}{exo:g11:trigo:1}
إلى \omterm{def:g11:trigo:radian}{الراديان}: $30^\circ = \frac{\pi}{6}$ و $45^\circ = \frac{\pi}{4}$ و
$120^\circ = \frac{2\pi}{3}$ و $270^\circ = \frac{3\pi}{2}$.

وإلى الدرجات: $\frac{\pi}{5} = 36^\circ$ و $\frac{5\pi}{6} = 150^\circ$ و
$\frac{7\pi}{4} = 315^\circ$ و $\frac{2\pi}{9} = 40^\circ$.
\end{solution}

\begin{solution}{exo:g11:trigo:2}
القيمتان $\frac{2\pi}{3}$ و $-\frac{\pi}{4}$ موجودتان أصلًا في
$\intoc{-\pi}{\pi}$: في الربع الثاني والربع الرابع على \omterm{def:g10:coordgeom:system}{الترتيب}.

$\frac{17\pi}{6} - 2\pi = \frac{5\pi}{6}$: في الربع الثاني، قريبًا من
المحور $x$ السالب.

$-\frac{7\pi}{2} + 4\pi = \frac{\pi}{2}$: أعلى نقطة في الدائرة،
$(0, 1)$.
\end{solution}

\begin{solution}{exo:g11:trigo:3}
$\cos\frac{3\pi}{4} = \cos\left(\pi - \frac\pi4\right)
= -\cos\frac\pi4 = -\frac{\sqrt2}{2}$.

$\sin\frac{5\pi}{6} = \sin\left(\pi - \frac\pi6\right)
= \sin\frac\pi6 = \frac12$.

$\cos\left(-\frac\pi3\right) = \cos\frac\pi3 = \frac12$.

$\sin\frac{4\pi}{3} = \sin\left(\pi + \frac\pi3\right)
= -\sin\frac\pi3 = -\frac{\sqrt3}{2}$.

$\cos\frac{11\pi}{6} = \cos\left(-\frac\pi6 + 2\pi\right)
= \cos\frac\pi6 = \frac{\sqrt3}{2}$.
\end{solution}

\begin{solution}{exo:g11:trigo:4}
$\sin^2 t = 1 - \cos^2 t = 1 - \frac{25}{169} = \frac{144}{169}$، إذن
$\sin t = \pm\frac{12}{13}$. وعلى $\intoo{-\frac\pi2}{0}$ (الربع الرابع)
يكون \omterm{def:g10:coordgeom:system}{الترتيب} سالبًا: $\sin t = -\frac{12}{13}$.
\end{solution}

\begin{solution}{exo:g11:trigo:5}
مستعملين \cref{prop:g11:trigo:associated}:
\[
A = (-\cos x) + \cos x + \cos x - \cos x = 0 ,
\]
\[
B = \sin x + (-\sin x) + (-\sin x) = -\sin x .
\]
\end{solution}

\begin{solution}{exo:g11:trigo:6}
$\cos t = \frac{\sqrt3}{2}$: أحد الحلول $\frac\pi6$، إذن
$t = \pm\frac\pi6$ (فالانسحابات بمقدار $2k\pi$ تخرج من \omterm{def:g10:numbers:interval}{الفترة}).

$\sin t = \frac12$: $t = \frac\pi6$ أو $t = \pi - \frac\pi6 =
\frac{5\pi}{6}$.

$\cos t = -1$: النقطة الوحيدة ذات \omterm{def:g10:coordgeom:system}{الفاصلة} $-1$ هي $M(\pi)$، إذن $t = \pi$.
\end{solution}

\begin{solution}{exo:g11:trigo:7}
$\sin t = -\frac{\sqrt3}{2}$: انطلاقًا من $a = -\frac\pi3$، تكون العائلتان
$-\frac\pi3 + 2k\pi$ و $\pi + \frac\pi3 + 2k\pi$؛ وفي
$\intco{0}{2\pi}$: $t = \frac{4\pi}{3}$ و $t = \frac{5\pi}{3}$.

$\cos t = 0$: $t = \frac\pi2$ و $t = \frac{3\pi}{2}$.

و $2\sin t + 1 = 0$ يعني $\sin t = -\frac12$: انطلاقًا من $a = -\frac\pi6$، في
$\intco{0}{2\pi}$: $t = \frac{7\pi}{6}$ و $t = \frac{11\pi}{6}$.
\end{solution}

\begin{solution}{exo:g11:trigo:8}
حسب \cref{thm:g11:trigo:equations}، تعني $\cos 2t = \cos t$ أن
$2t = t + 2k\pi$ أو $2t = -t + 2k\pi$، أي $t = 2k\pi$ أو
$t = \frac{2k\pi}{3}$ (حيث $k \in \Z$). والعائلة الأولى محتواة في
الثانية (بأخذ $k$ من مضاعفات $3$)، إذن مجموعة الحلول هي
\[
S = \left\{\frac{2k\pi}{3} : k \in \Z\right\}
\]
--- أي، على الدائرة، النقط الثلاث $M(0)$ و
$M\!\left(\frac{2\pi}{3}\right)$ و $M\!\left(\frac{4\pi}{3}\right)$.
\end{solution}

\begin{solution}{exo:g11:trigo:9}
\omterm{def:g10:algebra:expand}{بنشر} المربعين واستعمال $\cos^2 x + \sin^2 x = 1$:
\[
(\cos^2 x + 2\sin x\cos x + \sin^2 x)
+ (\cos^2 x - 2\sin x\cos x + \sin^2 x)
= 1 + 1 = 2 .
\]
\end{solution}

\begin{solution}{exo:g11:trigo:10}
التدحرج دون انزلاق يعني أن طول القوس يساوي المسافة
المقطوعة: فيعطي $r\theta = 150$ سم أن
$\theta = \frac{150}{30} = 5$ راد $= 5 \times \frac{180}{\pi} \approx
286.5^\circ$. وتقدّم الدورة الكاملة الدراجةَ بمقدار المحيط
$2\pi \times 30 = 60\pi \approx 188.5$ سم، أي نحو $1.88$ م.
\end{solution}

\begin{solution}{exo:g11:trigo:11}
على دائرة الوحدة، النقطتان اللتان فاصلتهما $\frac12$ بالضبط هما
$M\!\left(\pm\frac\pi3\right)$. والنقط التي \omterm{def:g10:coordgeom:system}{فاصلتها} \emph{لا تتجاوز}
$\frac12$ تشكّل القوس الواقع يسار المستقيم الشاقولي $x = \frac12$،
المقطوع عندما يسير $t$ من $\frac\pi3$ إلى $\pi$، أو من $-\pi$ إلى
$-\frac\pi3$. ومنه، على $\intoc{-\pi}{\pi}$:
\[
S = \intoc{-\pi}{-\frac{\pi}{3}} \cup \intcc{\frac{\pi}{3}}{\pi}.
\]
(والطرفان $\pm\frac\pi3$ مُدرَجان لأن المتراجحة واسعة.)
\end{solution}

\begin{solution}{pb:g11:trigo:1}
\textbf{1.} $30^\circ = \frac{\pi}{6}$؛ $135^\circ =
\frac{3\pi}{4}$؛ $300^\circ = \frac{5\pi}{3}$. و
$\frac{3\pi}{4} = 135^\circ$؛ $\frac{5\pi}{6} = 150^\circ$.

\textbf{2.} الطول $r\,t = 3 \times 2 = 6$ م. \omterm{def:g11:trigo:radian}{فالراديان} يجعل طول
القوس جداءً بسيطًا --- أما الدرجات فتجرّ
$\frac{\pi}{180}$ إلى كل صيغة.

\textbf{3.} عدد الراديانات المقطوعة $= \frac{\text{المسافة}}
{\text{نصف القطر}} = \frac{1000}{0.30} \approx 3\,333$ راد؛
وبالقسمة على $2\pi$: نحو $530$ دورة كاملة.

\textbf{4.} $\cos\frac{2\pi}{3} = -\frac12$؛
و $\sin\left(-\frac\pi4\right) = -\frac{\sqrt2}{2}$؛
و $\frac{19\pi}{6} - 2\pi = \frac{7\pi}{6}$، إذن
$\cos\frac{19\pi}{6} = \cos\frac{7\pi}{6} =
-\frac{\sqrt3}{2}$.

\textbf{5.} $\cos^2 t = 1 - \frac{9}{25} = \frac{16}{25}$؛ وفي
الربع الثاني يكون \omterm{def:g11:trigo:cossin}{جيب التمام} سالبًا:
$\cos t = -\frac45$، و
$\tan t = \frac{3/5}{-4/5} = -\frac34$.

\textbf{6.} في $t$ دقيقة تدور العجلة
$\frac{t}{30} \times 2\pi = \frac{\pi t}{15}$. وعند $t = 0$:
$h = 65 - 60 = 5$ م --- وهي منصة الصعود في الأسفل
(ارتفاع المحور ناقص نصف القطر). وإشارة الطرح تضعك \emph{تحت}
المحور عندما تكون الزاوية المقطوعة $0$.

\textbf{7.} $h(7.5) = 65 - 60\cos\frac{\pi}{2} = 65$ م (ارتفاع
المحور)؛ و $h(15) = 65 + 60 = 125$ م (القمة)؛
و $h(20) = 65 - 60\cos\frac{4\pi}{3} = 65 + 30 = 95$ م.

\textbf{8.} يعطي $65 - 60\cos\frac{\pi t}{15} = 95$ أن
$\cos\frac{\pi t}{15} = -\frac12$: فالحل الأول
$\frac{\pi t}{15} = \frac{2\pi}{3}$، أي $t = 10$ دقائق.

\textbf{9.} $\cos\frac{\pi t}{15} \leq -\frac12$ من أجل
$\frac{\pi t}{15} \in \intcc{\frac{2\pi}{3}}{\frac{4\pi}{3}}$،
أي $t \in \intcc{10}{20}$: أي عشر دقائق من الجولة فوق
$95$ م، متناظرة حول القمة عند $t = 15$.

\textbf{10.} الدقيقة الأولى:
\[
h(1) - h(0) = 60\left(1 - \cos\frac{\pi}{15}\right)
\approx 1.3 \text{ م}.
\]
وبين الدقيقتين $7$ و $8$:
\[
60\left(\cos\frac{7\pi}{15} - \cos\frac{8\pi}{15}\right)
\approx 12.5 \text{ م}
\]
--- أي نحو عشرة أضعاف. فالارتفاع يتغير أسرع ما يكون عند المرور بمستوى المحور
ويتباطأ قرب الأسفل والقمة: فالجيبية شديدة الانحدار في
وسطها، مسطحة عند طرفيها.

\textbf{11.} نصف القطر $= \frac{40 - 4}{2} = 18$ م، والمحور عند
$4 + 18 = 22$ م، والدور $12$ دقيقة:
$h(t) = 22 - 18\cos\left(\frac{\pi t}{6}\right)$.

\textbf{12.} \omterm{def:g10:functions:extrema}{القيمة العظمى} $7 + 3 = 10$ م عندما
$\sin\frac{\pi t}{6} = 1$: أول مرة عند $t = 3$ (الساعة 3 صباحًا)؛
\omterm{def:g10:functions:extrema}{والقيمة الصغرى} $4$ م أول مرة عند $t = 9$؛ والدور
$\frac{2\pi}{\pi/6} = 12$ ساعة.

\textbf{13.} $\sin\frac{\pi t}{6} \geq \frac12$ من أجل
$\frac{\pi t}{6} \in \intcc{\frac\pi6}{\frac{5\pi}{6}}$:
$t \in \intcc{1}{5}$ --- أي نافذة من أربع ساعات من الساعة 1 صباحًا إلى 5
صباحًا، تُفتح من جديد بعد دور واحد، $t \in \intcc{13}{17}$
(من 1 ظهرًا إلى 5 عصرًا).

\textbf{14.} المقابلة بالنسبة إلى المحور الشاقولي ترسل النقطة ذات
الزاوية $t$ إلى النقطة ذات الزاوية $\pi - t$، محافظةً على
\omterm{def:g10:coordgeom:system}{الترتيب}: $\sin(\pi - t) = \sin t$. ونصف الدورة حول
المركز يرسل $t$ إلى $\pi + t$، فيغيّر إشارتي الإحداثيين:
$\cos(\pi + t) = -\cos t$. والمقابلة بالنسبة إلى القطر
تبدّل \omterm{def:g10:coordgeom:system}{الفاصلة} \omterm{def:g10:coordgeom:system}{والترتيب}:
$\cos\left(\frac\pi2 - t\right) = \sin t$ --- وهي
``التمام'' في الزوايا المتتامة.

\textbf{15.} $\sin(-t) = -\sin t$: \omterm{def:g11:func:parity}{فردية} (بتناظر منحنيها
بنصف دورة)؛ و $\cos(-t) = \cos t$: \omterm{def:g11:func:parity}{زوجية} (بالمرآة). و
$\sin t + t^3$ مجموع دالتين فرديتين: فهي \omterm{def:g11:func:parity}{فردية}.

\textbf{16.} \omterm{def:g10:algebra:expand}{بالتحليل إلى عوامل}: $(2\sin t + 1)(\sin t - 1) = 0$:
\omterm{def:g10:algebra:equation}{فالمعادلة} $\sin t = 1$ تعطي $t = \frac\pi2$؛ و
$\sin t = -\frac12$ يعطي $t = \frac{7\pi}{6}$ و
$t = \frac{11\pi}{6}$. إذن ثلاثة حلول على الدائرة.

\textbf{17.} \omterm{def:g11:trigo:radian}{الراديان} الواحد $= \frac{180}{\pi} \approx 57.3^\circ$.
إذن $\sin(1) \approx 0.841$ بينما $\sin(1^\circ) \approx
0.0175$: أي بمعامل يقارب $50$. والمغزى: تحقق من ضوء النمط
قبل الوثوق بأي عرض مثلثي.

\textbf{18.} تتساوى \omterm{def:g10:coordgeom:system}{الفاصلة} \omterm{def:g10:coordgeom:system}{والترتيب} على القطر:
$t = \frac\pi4$ و $t = \frac{5\pi}{4}$.

\textbf{19.} $\sin(0.1) = 0.09983\ldots$ مقابل $0.1$:
فالخطأ النسبي نحو $0.17\,\%$. فمن أجل الزوايا الصغيرة يعانق الوتر
القوس --- والنهاية في السنة الأخيرة
$\frac{\sin t}{t} \to 1$ هي بالضبط هذه الملاحظة وقد صارت
مبرهنة.

\textbf{20.} \omterm{def:g11:trigo:radian}{الراديان}: الزاوية $=$ القوس على دائرة الوحدة، فتُحسب
الزوايا كما تُحسب الأطوال. \omterm{def:g11:trigo:cossin}{وجيب التمام} \omterm{def:g11:trigo:cossin}{والجيب}: إحداثيا
عقرب الساعة --- ظله على الأرض وارتفاعه على
الجدار. وتناظرات الدائرة: كتالوغ المتطابقات كله،
مقروءًا على المرايا وأنصاف الدورات. ومعادلات
$\cos t = c$: جدول مواعيد الساعة، والدائرة تعرض
كل الحلول دفعة واحدة. والأثواب: ارتدت العجلة
$65 - 60\cos$، وارتدى المدّ $7 + 3\sin$ --- الساعة نفسها والرياضيات
نفسها.
\end{solution}
